The experiments address four questions: whether the stored cuts are valid in
practice, and what certification prevents; whether the cuts change SCIP's
results on benchmark models, including models that SCIP finds hard; what the
separator does and what it costs; and how much of the root gap the cuts close
on instances with the structure of Section~\ref{sec:composition}, compared
with SCIP's own alternatives and with Gurobi. We report four campaigns. Each
was run with code, limits and protocol fixed before its first run; we call
these results \emph{prospective}, and we label analyses designed after
results were seen as \emph{post hoc}. Campaigns 1 and~2 used earlier
implementations and are summarized in Appendix~\ref{app:campaigns12}.
Campaign~3 tests the final implementation on MINLPLib models (Parts~A and~B)
and on a constructed family (Part~C). Campaign~4 was designed after campaign~3
to answer the questions that its results raised: it adds comparators, fresh
instances, larger models, a variant of the family in which the block minima
do not give the optimum, and an ablation of the certificate. Table~\ref{tab:overview} lists the parts.

\subsection{Setup}\label{sec:setup}

All runs use SCIP~10.0.2 through PySCIPOpt~6.2.1 and Python~3.13, with one
SCIP thread and one BLAS thread per process, on a Linux host (WSL2) with 18
cores and 36 hardware threads that was shared with unrelated jobs; the
one-minute load average was 14 to 21 during the timed runs of campaign~3 and
about 3 to 8 during campaign~4. At most six runs of a campaign were active at a
time; the modes of one instance run back to back in
rotated order, and we compare modes only within such paired runs. Timings are
therefore descriptive. Times are in-process wall times charged to the soft
budget (model build, discovery, separation and solve), excluding interpreter
start-up and the primal check; \emph{SCIP time excluding the callback} is
SCIP's solving time minus the time spent in the separator callback. Gurobi
13.0.3 \citep{Gurobi13} solves the same model with one thread, its nonconvex
mode (\code{NonConvex=2}) and the same gap limit. Models
come from MINLPLib \citep{BussieckDrudMeeraus2003,Vigerske2026MINLPLib} in
OSiL format (Appendix~\ref{app:instances}), stratified by the library's
metadata into convex models, nonconvex models without integer variables, and
nonconvex models with integer variables.

A model counts as \emph{solved} if SCIP (or Gurobi) reports optimality or the
relative gap limit $10^{-4}$, and the returned incumbent passes an independent
check of all original rows, bounds, domains, integrality restrictions and the
objective at a scaled tolerance of $10^{-5}$; the checker evaluates the source
expressions without the model builder. Times are compared by shifted geometric
means (shift 1\,s) and by the median of per-run ratios, over runs solved in
every compared mode. Dual bounds $a$ and $b$ are compared with the tolerance
$10^{-4}\max\{1,\abs a,\abs b\}$, the gap limit. The \emph{root gap closed} by a
mode is $(z_{\text{mode}}-z_{\text{base}})/(z^\ast-z_{\text{base}})$ for
minimization, where $z$ denotes root dual bounds and $z^\ast$ the optimal or
best known value; for a run that ends at the root without reporting a root
bound, because the root node was pruned, $z$ is its final dual bound. Every recorded cut is replayed in a fresh process
(Section~\ref{sec:replay}). A run whose worker failed has no complete cut log;
its cut count is reported as unknown, not as zero.

\begin{table}[t]
\centering\footnotesize
\caption{Parts of campaigns 3 and 4. Full runs have a budget of 300\,s; root
runs have a node limit of one. All parts except 3D are prospective. Modes
are defined in Appendix~\ref{app:software} (Table~\ref{tab:modes}); ``extra''
is SCIP with its disabled nonconvex separators switched on.}\label{tab:overview}
\begin{tabular}{@{}l>{\raggedright\arraybackslash}p{0.33\textwidth}>{\raggedright\arraybackslash}p{0.43\textwidth}l@{}}
\toprule
Part & Instances & Modes & Runs\\
\midrule
3A & 30 MINLPLib models of campaign 2, seeds 0--2 & baseline, all, auto; root runs also all-diag & full, root\\
3B & 30 structure-selected models, seeds 0--1 & baseline, all, auto; root runs also all-diag & full, root\\
3C & 20 path-family instances & baseline, all-diag-mech & full, root\\
3D & as 3A--3C (post hoc) & whole-row variant & as 3A--3C\\
\midrule
4C2 & the 20 instances of 3C & checkvarlocks off, extra, frozen-wide, Gurobi & full, root\\
4C3 & 20 fresh path-family instances & baseline, all-diag-mech, rowdir-wide, extra, Gurobi & full, root\\
4C4 & the instances of 4C3 with a binding row & baseline, frozen-wide, rowdir-wide, extra, Gurobi & full, root\\
4B2 & the 30 models of 3B & no aggregation: baseline, all, all-diag, all-diag-rowdir; extra & root\\
4D & 20 larger models that SCIP finds hard & baseline, all, auto, extra; root runs also all-diag, all-diag-rowdir & full, root\\
4U & recorded cuts of campaigns 3 and 4 & uncertified constants & offline\\
4S & random stars with $10$ to $10^4$ leaves & two star algorithms & offline\\
\bottomrule
\end{tabular}
\end{table}
